% !TeX program = xelatex
% 计算方法 作业01 —— LaTeX 排版
% 编译：XeLaTeX（中文需 ctex 宏包，建议 Overleaf 或 TeX Live 完整版）
\documentclass[a4paper,11pt]{ctexart}

\usepackage[margin=2.5cm]{geometry}
\usepackage{amsmath, amssymb}
\usepackage{booktabs}
\usepackage{array}
\usepackage{xcolor}
\usepackage{listings}

% ---------- MATLAB 代码样式 ----------
\definecolor{codegreen}{rgb}{0.1,0.55,0.1}
\definecolor{codegray}{rgb}{0.45,0.45,0.45}
\definecolor{codeblue}{rgb}{0.0,0.0,0.75}
\definecolor{codebg}{rgb}{0.97,0.97,0.97}
\lstset{
  language=Matlab,
  basicstyle=\ttfamily\small,
  commentstyle=\color{codegreen},
  keywordstyle=\color{codeblue},
  stringstyle=\color{red!70!black},
  numbers=left,
  numberstyle=\tiny\color{codegray},
  stepnumber=1,
  numbersep=5pt,
  backgroundcolor=\color{codebg},
  frame=single,
  rulecolor=\color{black!25},
  breaklines=true,
  showstringspaces=false,
  tabsize=4,
  extendedchars=true,
}

\setlength{\parindent}{0pt}   % 中文习惯：段落不缩进
\setlength{\parskip}{0.5em}   % 段间距

\begin{document}

% =====================================================================
% 封面
% =====================================================================
\begin{titlepage}
  \centering
  {\zihao{2}\bfseries 深圳大学医学部生物医学工程学院\par}
  \vspace{1.2em}
  {\zihao{1}\bfseries 本科生课程作业\par}
  \vspace{2.5em}
  {\zihao{4}
    课程：计算方法\\
    （2026--2027学年第一学期）\\
    任课教师：张力\par}
  \vspace{2.5em}

  {\zihao{4}
  \renewcommand{\arraystretch}{1.6}
  \begin{tabular}{p{3.5em}p{18em}}
    班级 & 计算方法\_\_\_\_\_\_班 \\[0.4em]
    学号 & 2024220029 \\[0.4em]
    姓名 & \\[0.4em]
    序号 & \\[0.4em]
    提交日期 & \_\_\_\_年 \_\_\_\_ 月 \_\_\_\_ 日 \\
  \end{tabular}
  \par}

  \vspace{2.5em}
  {\zihao{4}\bfseries 供助教评分使用\par}
  \vspace{1em}
  {\zihao{4}
  \renewcommand{\arraystretch}{1.6}
  \begin{tabular}{p{4.5em}p{17em}}
    助教姓名 & \\[0.4em]
    收到日期 & \_\_\_\_年 \_\_\_\_ 月 \_\_\_\_ 日 \\[0.4em]
    评分(0--100) & \\[0.4em]
    评语（如有） & \\[0.4em]
  \end{tabular}
  \par}
\end{titlepage}

% =====================================================================
% 作业正文
% =====================================================================
\setcounter{section}{0}

% ---------- 第 1 题 ----------
\section*{1.\ 有效数字位数（20分）}
\addcontentsline{toc}{section}{1.\ 有效数字位数}

\textbf{题目}　设准确值为 $x=3.78695$，$y=10$，它们的近似值分别为
$x_1^*=3.7869$，$x_2^*=3.7870$ 及 $y_1^*=9.9999$，$y_2^*=10.1$，$y_3^*=10.0001$，
试分析 $x_1^*,\,x_2^*,\,y_1^*,\,y_2^*,\,y_3^*$ 分别具有几位有效数字。

\textbf{解}

判定准则：令准确值 $x$ 的近似值为 $x^*$，将 $x^*$ 写成
$x^*=0.a_1a_2\cdots\times 10^{m}$（$a_1\neq 0$）的规格化形式，若绝对误差
$|x^*-x|\le 0.5\times 10^{m-n}$，则 $x^*$ 具有 $n$ 位有效数字。

\begin{enumerate}
  \item $x_1^*=3.7869$，准确值 $x=3.78695$：
  \[
    |x_1^*-x|=|3.7869-3.78695|=0.00005=0.5\times 10^{-4}.
  \]
  $3.7869=0.37869\times 10^{1}$，即 $m=1$，由 $0.5\times 10^{1-n}=0.5\times 10^{-4}$
  得 $n=5$。故 $x_1^*$ 有 $5$ 位有效数字。

  \item $x_2^*=3.7870$，准确值 $x=3.78695$：
  \[
    |x_2^*-x|=0.00005=0.5\times 10^{-4},
  \]
  同 (1)，$n=5$。故 $x_2^*$ 有 $5$ 位有效数字。

  \item $y_1^*=9.9999$，准确值 $y=10$：
  \[
    |y_1^*-y|=0.0001=10^{-4}.
  \]
  $9.9999=0.99999\times 10^{1}$，即 $m=1$。
  $0.5\times 10^{1-4}=0.5\times 10^{-3}=0.0005\ge 0.0001$，而
  $0.5\times 10^{1-5}=0.00005<0.0001$。故 $y_1^*$ 有 $4$ 位有效数字。

  \item $y_2^*=10.1$，准确值 $y=10$：
  \[
    |y_2^*-y|=0.1=10^{-1}.
  \]
  $10.1=0.101\times 10^{2}$，即 $m=2$。
  $0.5\times 10^{2-2}=0.5\ge 0.1$，而 $0.5\times 10^{2-3}=0.05<0.1$。
  故 $y_2^*$ 有 $2$ 位有效数字。

  \item $y_3^*=10.0001$，准确值 $y=10$：
  \[
    |y_3^*-y|=0.0001=10^{-4}.
  \]
  $10.0001=0.100001\times 10^{2}$，即 $m=2$。
  $0.5\times 10^{2-5}=0.5\times 10^{-3}=0.0005\ge 0.0001$，而
  $0.5\times 10^{2-6}=0.00005<0.0001$。故 $y_3^*$ 有 $5$ 位有效数字。
\end{enumerate}

综上：
\begin{center}
  \begin{tabular}{lccccc}
    \toprule
    近似值 & $x_1^*$ & $x_2^*$ & $y_1^*$ & $y_2^*$ & $y_3^*$ \\
    \midrule
    有效数字位数 & 5 位 & 5 位 & 4 位 & 2 位 & 5 位 \\
    \bottomrule
  \end{tabular}
\end{center}

% ---------- 第 2 题 ----------
\section*{2.\ Gauss 消去法（20分）}
\addcontentsline{toc}{section}{2.\ Gauss 消去法}

\textbf{题目}　用（1）矩阵左乘实现 Gauss 消去法，与（2）Gauss 列主元素消去法
求解下列方程组 $Ax=b$，其中
\[
  A=\begin{pmatrix}
    1 & 2 & 1 & -2\\
    2 & 5 & 3 & -2\\
    -2 & -2 & 3 & 5\\
    1 & 3 & 2 & 5
  \end{pmatrix},\qquad
  b=\begin{pmatrix}-1\\3\\15\\9\end{pmatrix}.
\]

\textbf{解}

增广矩阵 $[A\mid b]$：
\[
  \left(\begin{array}{rrrr|r}
    1 & 2 & 1 & -2 & -1\\
    2 & 5 & 3 & -2 & 3\\
    -2 & -2 & 3 & 5 & 15\\
    1 & 3 & 2 & 5 & 9
  \end{array}\right)
\]

\subsection*{（1）矩阵左乘实现 Gauss 消去法（顺序消元）}

\textbf{第 1 步}　主元 $a_{11}=1$，消去第 1 列下方元素
（行2 $-2\times$行1，行3 $+2\times$行1，行4 $-$行1）：
\[
  \left(\begin{array}{rrrr|r}
    1 & 2 & 1 & -2 & -1\\
    0 & 1 & 1 & 2 & 5\\
    0 & 2 & 5 & 1 & 13\\
    0 & 1 & 1 & 7 & 10
  \end{array}\right)
\]

\textbf{第 2 步}　主元 $a_{22}=1$，消去第 2 列下方元素
（行3 $-2\times$行2，行4 $-$行2）：
\[
  \left(\begin{array}{rrrr|r}
    1 & 2 & 1 & -2 & -1\\
    0 & 1 & 1 & 2 & 5\\
    0 & 0 & 3 & -3 & 3\\
    0 & 0 & 0 & 5 & 5
  \end{array}\right)
\]

\textbf{第 3 步}　主元 $a_{33}=3$，其下方已为 $0$，无需消元。此时系数矩阵已化为上三角矩阵。

\textbf{回代：}
\[
  \begin{aligned}
    5x_4&=5 & &\Rightarrow & x_4&=1\\
    3x_3-3x_4&=3 & &\Rightarrow & x_3&=2\\
    x_2+x_3+2x_4&=5 & &\Rightarrow & x_2&=1\\
    x_1+2x_2+x_3-2x_4&=-1 & &\Rightarrow & x_1&=-3
  \end{aligned}
\]
故解为 $x=(x_1,x_2,x_3,x_4)=(-3,1,2,1)$。

\subsection*{（2）Gauss 列主元素消去法}

\textbf{第 1 步}　第 1 列元素为 $1,2,-2,1$，绝对值最大者为 $2$（第 2 行），交换第 1、2 行：
\[
  \left(\begin{array}{rrrr|r}
    2 & 5 & 3 & -2 & 3\\
    1 & 2 & 1 & -2 & -1\\
    -2 & -2 & 3 & 5 & 15\\
    1 & 3 & 2 & 5 & 9
  \end{array}\right)
\]
以 $a_{11}=2$ 为主元消元（行2 $-\tfrac12\times$行1，行3 $+$行1，行4 $-\tfrac12\times$行1）：
\[
  \left(\begin{array}{rrrr|r}
    2 & 5 & 3 & -2 & 3\\
    0 & -\tfrac12 & -\tfrac12 & -1 & -\tfrac52\\
    0 & 3 & 6 & 3 & 18\\
    0 & \tfrac12 & \tfrac12 & 6 & \tfrac{15}{2}
  \end{array}\right)
\]

\textbf{第 2 步}　第 2 列（第 2 行起）元素为 $-\tfrac12,\,3,\,\tfrac12$，绝对值最大者为 $3$
（第 3 行），交换第 2、3 行：
\[
  \left(\begin{array}{rrrr|r}
    2 & 5 & 3 & -2 & 3\\
    0 & 3 & 6 & 3 & 18\\
    0 & -\tfrac12 & -\tfrac12 & -1 & -\tfrac52\\
    0 & \tfrac12 & \tfrac12 & 6 & \tfrac{15}{2}
  \end{array}\right)
\]
以 $a_{22}=3$ 为主元消元（行3 $+\tfrac16\times$行2，行4 $-\tfrac16\times$行2）：
\[
  \left(\begin{array}{rrrr|r}
    2 & 5 & 3 & -2 & 3\\
    0 & 3 & 6 & 3 & 18\\
    0 & 0 & \tfrac12 & -\tfrac12 & \tfrac12\\
    0 & 0 & -\tfrac12 & \tfrac{11}{2} & \tfrac92
  \end{array}\right)
\]

\textbf{第 3 步}　第 3 列（第 3 行起）元素为 $\tfrac12,\,-\tfrac12$，绝对值最大者取 $\tfrac12$
（第 3 行），不交换。以 $a_{33}=\tfrac12$ 为主元消元（行4 $+$行3）：
\[
  \left(\begin{array}{rrrr|r}
    2 & 5 & 3 & -2 & 3\\
    0 & 3 & 6 & 3 & 18\\
    0 & 0 & \tfrac12 & -\tfrac12 & \tfrac12\\
    0 & 0 & 0 & 5 & 5
  \end{array}\right)
\]

\textbf{回代：}
\[
  \begin{aligned}
    5x_4&=5 & &\Rightarrow & x_4&=1\\
    \tfrac12 x_3-\tfrac12 x_4&=\tfrac12 & &\Rightarrow & x_3&=2\\
    3x_2+6x_3+3x_4&=18 & &\Rightarrow & x_2&=1\\
    2x_1+5x_2+3x_3-2x_4&=3 & &\Rightarrow & x_1&=-3
  \end{aligned}
\]
综上，$x=(x_1,x_2,x_3,x_4)=(-3,1,2,1)$，与（1）结果一致。

% ---------- 第 3 题 ----------
\section*{3.\ Cholesky 分解（20分）}
\addcontentsline{toc}{section}{3.\ Cholesky 分解}

\textbf{题目}　求下列矩阵的 Cholesky 分解（可直接使用公式）：
\[
  A=\begin{pmatrix}
    1 & \tfrac12 & \tfrac13\\
    \tfrac12 & \tfrac13 & \tfrac14\\
    \tfrac13 & \tfrac14 & \tfrac15
  \end{pmatrix}.
\]

\textbf{解}

$A$ 为 $3\times3$ Hilbert 矩阵，对称正定。Cholesky 分解 $A=LL^{T}$，其中 $L$ 为
下三角矩阵（对角元为正）。递推公式（$i\ge j$）：
\[
  l_{ii}=\sqrt{a_{ii}-\sum_{k=1}^{i-1}l_{ik}^{2}},
  \qquad
  l_{ij}=\frac{a_{ij}-\sum_{k=1}^{j-1}l_{ik}l_{jk}}{l_{jj}}\quad(i>j).
\]

计算过程：
\[
  \begin{aligned}
    l_{11}&=\sqrt{a_{11}}=1,\\
    l_{21}&=\frac{a_{21}}{l_{11}}=\frac12,\qquad
    l_{31}=\frac{a_{31}}{l_{11}}=\frac13,\\
    l_{22}&=\sqrt{a_{22}-l_{21}^{2}}=\sqrt{\tfrac13-\tfrac14}=\sqrt{\tfrac{1}{12}}=\frac{1}{\sqrt{12}},\\
    l_{32}&=\frac{a_{32}-l_{31}l_{21}}{l_{22}}
          =\frac{\tfrac14-\tfrac16}{1/\sqrt{12}}=\frac{1}{\sqrt{12}},\\
    l_{33}&=\sqrt{a_{33}-l_{31}^{2}-l_{32}^{2}}
          =\sqrt{\tfrac15-\tfrac19-\tfrac{1}{12}}
          =\sqrt{\tfrac{1}{180}}=\frac{1}{\sqrt{180}}.
  \end{aligned}
\]

结果：
\[
  L=\begin{pmatrix}
    1 & 0 & 0\\
    \tfrac12 & \tfrac{1}{\sqrt{12}} & 0\\
    \tfrac13 & \tfrac{1}{\sqrt{12}} & \tfrac{1}{\sqrt{180}}
  \end{pmatrix}.
\]
（等价地，有理化后 $l_{22}=l_{32}=\dfrac{\sqrt3}{6}\approx0.2887$，
$l_{33}=\dfrac{\sqrt5}{30}\approx0.0745$。）

% ---------- 第 4 题 ----------
\section*{4.\ 追赶法（Crout 法）（20分）}
\addcontentsline{toc}{section}{4.\ 追赶法（Crout 法）}

\textbf{题目}　用追赶法（Crout 法）解三对角方程组 $Ax=f$（可直接使用公式），其中
\[
  A=\begin{pmatrix}
    2 & 1 & 0 & 0\\
    1 & 4 & 1 & 0\\
    0 & 1 & 4 & 1\\
    0 & 0 & 1 & 2
  \end{pmatrix},\qquad
  f=\begin{pmatrix}-3\\6\\14\\-2\end{pmatrix}.
\]

\textbf{解}

记主对角线 $b=(2,4,4,2)$，下对角线 $a=(1,1,1)$，上对角线 $c=(1,1,1)$。
追赶法（Crout 分解）将 $A$ 分解为 $A=LU$：$L$ 为下三角矩阵（主对角元素记为 $p_i$，
副对角元素为 $a_i$），$U$ 为单位上三角矩阵（副对角元素记为 $q_i$）。

递推公式：
\[
  p_1=b_1,\qquad q_i=\frac{c_i}{p_i},\qquad p_i=b_i-a_iq_{i-1}\quad(i=2,3,4).
\]

计算 $L$、$U$ 的系数：
\[
  \begin{aligned}
    p_1&=2, & q_1&=\frac{c_1}{p_1}=\frac12,\\
    p_2&=b_2-a_2q_1=4-\tfrac12=\tfrac72, & q_2&=\frac{c_2}{p_2}=\frac27,\\
    p_3&=b_3-a_3q_2=4-\tfrac27=\tfrac{26}{7}, & q_3&=\frac{c_3}{p_3}=\frac{7}{26},\\
    p_4&=b_4-a_4q_3=2-\tfrac{7}{26}=\tfrac{45}{26}.
  \end{aligned}
\]
即 $p=(2,\tfrac72,\tfrac{26}{7},\tfrac{45}{26})$，$q=(\tfrac12,\tfrac27,\tfrac{7}{26})$。

``追''——前推解 $Ly=f$：
\[
  \begin{aligned}
    2y_1&=-3 & &\Rightarrow & y_1&=-\tfrac32\\
    y_1+\tfrac72 y_2&=6 & &\Rightarrow & y_2&=\tfrac{15}{7}\\
    y_2+\tfrac{26}{7}y_3&=14 & &\Rightarrow & y_3&=\tfrac{83}{26}\\
    y_3+\tfrac{45}{26}y_4&=-2 & &\Rightarrow & y_4&=-3
  \end{aligned}
\]

``赶''——回代解 $Ux=y$：
\[
  \begin{aligned}
    x_4&=y_4=-3\\
    x_3+\tfrac{7}{26}x_4&=\tfrac{83}{26} & &\Rightarrow & x_3&=4\\
    x_2+\tfrac27 x_3&=\tfrac{15}{7} & &\Rightarrow & x_2&=1\\
    x_1+\tfrac12 x_2&=-\tfrac32 & &\Rightarrow & x_1&=-2
  \end{aligned}
\]

综上：$x=(x_1,x_2,x_3,x_4)=(-2,1,4,-3)$。

% ---------- 第 5 题 ----------
\section*{5.\ MATLAB 程序 myDoolittle（20分）}
\addcontentsline{toc}{section}{5.\ MATLAB 程序 myDoolittle}

\textbf{题目}　使用 MATLAB 编写程序 \texttt{myDoolittle}，要求使用 Doolittle 法实现 LU 分解，
使得程序输入矩阵 $A$，返回单位下三角矩阵 $L$ 和上三角矩阵 $U$。测试矩阵
$A=\begin{pmatrix}1&-1&3\\2&-4&6\\4&-9&2\end{pmatrix}$ 作为输入的输出结果是否正确。

\textbf{解}

完整程序：
\begin{lstlisting}
function [L, U] = myDoolittle(A)
% myDoolittle  用 Doolittle 法实现矩阵 A 的 LU 分解
%   输入：A —— n×n 可逆矩阵
%   输出：L —— 单位下三角矩阵（对角元全为 1）
%         U —— 上三角矩阵
%   满足：A = L*U
n = size(A, 1);
L = eye(n);     % 单位下三角矩阵
U = zeros(n);   % 上三角矩阵

for k = 1:n
    % (1) 先求 U 的第 k 行
    for j = k:n
        U(k, j) = A(k, j) - L(k, 1:k-1) * U(1:k-1, j);
    end
    % (2) 再求 L 的第 k 列
    for i = k+1:n
        L(i, k) = (A(i, k) - L(i, 1:k-1) * U(1:k-1, k)) / U(k, k);
    end
end
end
\end{lstlisting}

测试脚本：
\begin{lstlisting}
A = [1 -1 3; 2 -4 6; 4 -9 2];
[L, U] = myDoolittle(A);
disp(L), disp(U), disp(L*U)
\end{lstlisting}

运行输出：
\begin{lstlisting}
L =
     1     0     0
     2     1     0
     4   2.5     1
U =
     1    -1     3
     0    -2     0
     0     0   -10
L*U =
     1    -1     3
     2    -4     6
     4    -9     2
\end{lstlisting}

可见 $L\cdot U=A$，程序输出正确。

\end{document}
